\documentclass[11pt]{article}
\usepackage{geometry}
\geometry{margin=1in}

\title{Similarity Renormalization Group & Decoupling}
\author{Maximilian Nidens}
\date{28.01.2026}

\begin{document}

\begin{abstract}
Quantum field theories experience divergencies due to local interactions, thus requiring cutoffs to 
produce finite results which is called regularization.
With that, our physical observables are then cutoff dependent. However, our physical observables must remain independet of 
the chosen regulator, because physics must be scale invariant.\\
A renormalization group provides transformations evolving the Hamiltonian by absorbing scale-dependent terms while preserving
the observables. 
In low-energy nuclear physics, the similarity renormalization group approach essentially softens
the underlying potential with continous unitary transformations by decoupling low-energy degrees of freedom from the high-energy
regime. This is typically achieved through driving the potential into a band diagonal form in momentum space, such that 
one can cut off the high-energy regime of the potential, leading to faster convergence in actual calculations.
\end{abstract}

\end{document}